\documentclass[11pt]{article}
\usepackage[a4paper,margin=26mm]{geometry}
\usepackage{amsmath,amssymb}
\setlength{\parindent}{0pt}
\setlength{\parskip}{7pt}
\emergencystretch=2em
\title{English auction revenue need not strictly exceed first-price revenue\\
\large Disproof of conjecture 00000008235}
\author{ziangni-sys}
\date{4 October 2026}
\begin{document}
\maketitle
The English conjecture asserts that the sign of the revenue difference
between English and first-price auctions is always positive. The Chinese
version likewise says that the sign is strictly positive. We refute this
universal assertion with equal equilibrium expected revenues. This does
not refute a nonnegative linkage comparison.

\textbf{Private values and auction conventions.}
There are two bidders, each with value $1$ for a single indivisible item.
Each private-value distribution is the point mass at $1$. These values
are independent and identically distributed: for arbitrary measurable
sets $A,B$, the probabilities of their individual and joint membership
events are zero or one and satisfy the independence product identity.
The conjecture imposes no nondegeneracy restriction; the point-mass
assumption is explicit in this example.

Both auctions have zero reserve. Bids and dropout thresholds are
nonnegative. Ties award the item to bidder $0$. Utility is value of the
received item minus payment, with no payment by a losing bidder.

The first-price auction awards to the highest bidder and charges that
bid. The ascending English auction uses dropout thresholds: the clock
starts at zero and stops at the first dropout, selling to the remaining
bidder. If both drop simultaneously, the stated tie convention applies.
For threshold vector $(b_0,b_1)$ the stopping price is
\[
 p_{\rm stop}=\min(b_0,b_1),
\]
and the winner is the bidder with the larger threshold. This is the actual
ascending-clock stopping rule: no bidder drops at a nonnegative price
$q<p_{\rm stop}$, while at $p_{\rm stop}$ at least one does. Thus these
payments are derived from the English mechanism.

\textbf{Equilibrium in the first-price auction.}
Take the bid profile $(1,1)$. Bidder $0$ receives the item and pays $1$;
bidder $1$ loses and pays zero. Both utilities are zero.
A unilateral bid $r<1$ loses, giving zero utility. A bid $r>1$ wins but
gives utility $1-r<0$. A bid $r=1$ gives zero utility, with either tie outcome.
Thus no nonnegative unilateral deviation improves utility, including
all boundary and tie cases. This is a Nash equilibrium and, because
the type distribution is degenerate, a Bayesian Nash equilibrium.

\newpage
\textbf{Equilibrium in the English auction.}
Take dropout thresholds $(1,1)$. The first dropout price is $1$; the
tie winner pays $1$, so both utilities again vanish.
A unilateral threshold $r<1$ drops first and loses, giving zero utility.
For $r>1$, the opponent drops at $1$ and the deviator wins at price $1$,
again giving zero utility. At $r=1$ the tie convention gives zero utility.
Thus all threshold deviations have utility zero or less, and the profile
is a Nash and Bayesian Nash equilibrium. In the two-bidder ascending
dropout model these are all unilateral dropout choices.

\textbf{Actual expected revenues.}
Each auction sells the item at price $1$, so its revenue is the sum of
the actual bidder payments, identically $1$ on the probability space.
Consequently
\[
 \mathbb E R_{\rm English}=1=\mathbb E R_{\rm first-price},
 \qquad
 \mathbb E R_{\rm English}-\mathbb E R_{\rm first-price}=0.
\]
The difference is not strictly positive. Independence, zero reserve,
the tie convention, and degeneracy of the distribution are all stated,
rather than implicitly adding hypotheses to the conjecture. The example
makes no separate assertion about Myerson's reserve formula or
asymptotic gains under positive affiliation.

\textbf{Formal correspondence.}
Lean constructs a genuine \texttt{PMF.pure} probability law on the
one-point sample space and its associated probability measure.
\texttt{independent\_values} proves Mathlib's \texttt{IndepFun} predicate
for the two real private values. Allocations, winning-bid payments,
English dropout payments, utilities, and unilateral updates are defined
as actual functions.

\texttt{actual\_first\_dropout} proves the displayed stopping rule for
every admissible threshold profile; \texttt{dropout\_unique} proves
its uniqueness. \texttt{winner\_maximizes} proves the allocation selects
a highest bid, and \texttt{allocation\_feasible} proves exactly one
item is allocated. \texttt{revenue\_is\_payments} proves actual total
payments equal the mechanism price.

The equilibrium predicates quantify over both bidders and every
nonnegative unilateral report. \texttt{first\_equilibrium} and
\texttt{english\_equilibrium} verify them, covering all tie cases.
Expected revenue is the actual Bochner integral of total payments
against the constructed probability measure.
\texttt{actual\_revenues} and \texttt{counterexample} prove equality
of these integrals and the failure of strict positivity.

\textbf{Reproduction.}
Use Lean 4.19.0 and public Mathlib revision
\begin{center}\small\texttt{c44e0c8ee63ca166450922a373c7409c5d26b00b}\end{center}
Run \texttt{lake build} inside \texttt{lean/}. Axiom audits are printed.
The validation notes record the final full build and PDF inspection.
\end{document}
